\chapter{L'azienda Eggon}
\label{chap:introduzione}

\section{Identità e contesto di Eggon}

Presentazione dello studio di sviluppo videoludico indipendente \textit{Eggon}: dimensioni del team, tipologia di prodotti realizzati, posizionamento nel mercato \textit{consumer} su piattaforme digitali quali \textit{Steam}, \textit{App Store}, \textit{Google Play} e \textit{Nintendo eShop}. Le informazioni verranno riportate direttamente sulla base di quanto osservato durante il periodo di stage, tenendo conto che il team responsabile dello sviluppo originario del videogioco non faceva più parte dell'azienda.

Eggon è una \textit{software house} italiana, fondata a fine 2015 da Gianpaolo Ferrarin e Fulvio Menegozzo nel loro ritorno a Padova dopo un'esperienza 
startupper negli Stati Uniti, e conta circa 30 persona che vi lavorano.
Oggi è una realtà che opera su tre settori chiave:
- Digital Health, che con l'app MioPEDIATRA ha connesso centinaia di migliaia di genitori con i pediatri, facilitandone le operazioni come prescrizioni, appuntamenti e tanto altro.
- IoT e Industria 5.0, con software che puntano sull'efficientamento produttivo ed energetico.
- Insurtech, con l'app MyWide che permette una gestione facilitata delle polizze assicurative.
Un quarto settore su cui hanno sperimento con qualche progetto è l'ambito videoludico che sarà l'interesse di questo elaborato.
In quest'ultimo ambito l'azienda ha sperimentato diversi tipi di progetti ludici, in modo a volte indipendente e altre in collaborazione con altri studi.
I progetti sviluppati sono i seguenti:
- Eleston, gioco da tavolo strategico in cui i giocatori assuomo il ruolo di maghi impegnati in una sfida per dimostrare la loro padronanza degli elementi.
- Diamonst, RP in realtà aumentata, per mobile, con un campagna per il giocatore singolo, duelli multiplayer e una modalità "Virual Pet".
- Fangold, gioco che mescola MMORP e TC, per pc, combinando elementi tipici dei giochi di carte con meccaniche da gioco di ruolo online.
- Lost in the Dungeon, l'ultimo gioco sviluppato, rilasciato a Marzo 2018. SI tratta di un gioco di carte che mescola meccaniche degli RP e meccaniche dei dungeon crawler, disponibile su pc e mobile.
Con l'attività di stage che questo documento ha l'obiettivo di descrivere e analizzare l'azienda ha deciso di riaprire un'altra volta questo suo piccolo paragrafo di storia, il quale verrà approfondito nel resto del documento.

\section{Tecnologie e strumenti adottati}

Panoramica del \textit{tech stack} utilizzato all'interno dello studio: motore grafico \textit{Unity} (versione~6), linguaggio di programmazione \textit{C\#}, sistema di versionamento \textit{Git} ospitato su \textit{Bitbucket} e integrazione con \textit{Steamworks.NET} per la distribuzione sulla piattaforma \textit{Steam}. Verrà inoltre descritta l'architettura del progetto, organizzata su due livelli: da un lato le librerie proprietarie condivise tra i diversi progetti sviluppati da Eggon, dall'altro il codice specifico del titolo \textit{LITD}. Particolare attenzione sarà dedicata al sistema di eventi interno, che rappresenta il principale meccanismo di comunicazione tra i diversi componenti dell'applicazione.

\section{Processi di sviluppo e organizzazione del lavoro}

Descrizione delle modalità operative adottate nello studio: sviluppo manuale con generazione delle \textit{build} direttamente dall'editor di \textit{Unity}, assenza di una pipeline di integrazione continua automatizzata, utilizzo di file condivisi per il tracciamento delle attività come strumento informale di pianificazione e lavoro in coppia tra gli stagisti sotto la supervisione del tutor aziendale. Verrà evidenziato come tali processi riflettano le esigenze e le caratteristiche organizzative tipiche di uno studio di sviluppo indipendente.

\section{Propensione all'innovazione}

Analisi delle strategie adottate da Eggon per perseguire i propri obiettivi di mercato: espansione verso nuove piattaforme hardware, come \textit{Nintendo Switch}, aggiornamento continuo del motore grafico quale strategia di modernizzazione tecnologica e integrazione con ecosistemi di distribuzione consolidati. Tale propensione all'innovazione costituisce il collegamento diretto con le motivazioni che hanno portato alla definizione del progetto di stage, oggetto di approfondimento nel capitolo successivo.